Polynomial functors can be developed in different settings,
including locally Cartesian closed categories \cite{gambino2013},
lex categories \cite{weber2015},
and $\pi$-clans \cite{hua2026}.
In \Cref{sec:exponentiable},
we saw that the theory of preclans
with an appropriate notion of preclan-exponentiability subsumes these various cases.
This chapter reproduces the theory of polynomial functors in this more general setting.

\section{Polynomial signatures}
From now on, we fix a preclan $(\CC,\RR)$ with a terminal object $1$.

\begin{definition}[Polynomial signature]
  A (univariate) \emph{polynomial signature} in $(\CC,\RR)$ consists of
  a morphism $f : E \to B$ such that
  $B \to 1$ is an $\RR$-map
  and $f$ is $\RR$-exponentiable.
\end{definition}

For simplicity, we only work with univariate polynomials.
This can be straightforwardly adapted to the multivariate case:
a multivariate polynomial signature would consist of three maps
\[\begin{tikzcd}
	I & E & B & J
	\arrow["i"', from=1-2, to=1-1]
	\arrow["f", from=1-2, to=1-3]
	\arrow["j", from=1-3, to=1-4]
\end{tikzcd}\]
such that $j$ is an $\RR$-map and $f$ is $\RR$-exponentiable.

\begin{definition}
  The \emph{category $\Poly_{(\CC,\RR)}$ of polynomial signatures}
  in $(\CC,\RR)$ has as objects
  polynomial signatures, and morphisms between objects
  $f : E \to B$ and $f' : E' \to B'$ as pairs
  \[(\phi_1 : B \to B', \phi^\# : f' \times_{B'} \phi_1 \to E)\]
  of morphisms in $\CC$ such that
  $f \circ \phi^\# = \pi_B : f' \times_{B'} \phi_1 \to B$.
  \[
  \begin{tikzcd}
    {f' \times_{B'} \phi_1} & E & B \\
    {E'} && {B'}
    \arrow["{\phi^\#}", from=1-1, to=1-2]
    \arrow["{\pi_B}", bend left = 50, from=1-1, to=1-3]
    \arrow["{\pi_{E'}}"', from=1-1, to=2-1]
    \arrow["\lrcorner"{anchor=center, pos=0.125}, draw=none, from=1-1, to=2-3]
    \arrow["f", from=1-2, to=1-3]
    \arrow["{\phi_1}", from=1-3, to=2-3]
    \arrow["{f'}"', from=2-1, to=2-3]
  \end{tikzcd}\]
  We write $\phi = (\phi_1, \phi^\#) : f \to f'$.
  The identity morphism on $(f : E \to B)$ is given by
  computing the pullback along an identity map
  \[
    \begin{tikzcd}
      {f \times_{B} \id_B} & E & B \\
      E && B
      \arrow["{{\iso}}"{description}, draw=none, from=1-1, to=1-2]
      \arrow["{{\pi_{E}}}"', from=1-1, to=2-1]
      \arrow["f", from=1-2, to=1-3]
      \arrow["{\id_E}"{description}, from=1-2, to=2-1]
      \arrow["\lrcorner"{anchor=center, pos=0.125}, draw=none, from=1-2, to=2-3]
      \arrow["{{\id_B}}", from=1-3, to=2-3]
      \arrow["{{f}}"', from=2-1, to=2-3]
    \end{tikzcd}
  \]
  Composition of $\phi : (f_0 : E_0 \to B_0) \to (f_1 : E_1 \to B_1)$
  and $\psi : (f_1 : E_1 \to B_1) \to (f_2 : E_2 \to B_2)$
  is given as follows
  \[
  \begin{tikzcd}
    {f_2 \times_{B_2} (\psi_1 \circ \phi_1)} & {f_1 \times_{B_1} \phi_1} & {E_0} & {B_0} \\
    {f_2 \times_{B_2} \psi_1} & {E_1} && {B_1} \\
    {E_2} &&& {B_2}
    \arrow["{\phi_1^* \psi^\#}", from=1-1, to=1-2]
    \arrow[from=1-1, to=2-1]
    \arrow["\lrcorner"{anchor=center, pos=0.125}, draw=none, from=1-1, to=2-2]
    \arrow["{\phi^\#}", from=1-2, to=1-3]
    \arrow[from=1-2, to=2-2]
    \arrow["\lrcorner"{anchor=center, pos=0.125}, draw=none, from=1-2, to=2-4]
    \arrow["{f_0}", from=1-3, to=1-4]
    \arrow["{\phi_1}", from=1-4, to=2-4]
    \arrow["{\psi^\#}", from=2-1, to=2-2]
    \arrow[from=2-1, to=3-1]
    \arrow["\lrcorner"{anchor=center, pos=0.125}, draw=none, from=2-1, to=3-4]
    \arrow["{f_1}", from=2-2, to=2-4]
    \arrow["{\psi_1}", from=2-4, to=3-4]
    \arrow["{f_2}", from=3-1, to=3-4]
  \end{tikzcd}
  \]
\end{definition}

On the notational difference with \cite{gambino2013}:
when $\CC$ is LCC (locally cartesian closed) and $\RR$ is the class of all maps,
our 1-category $\Poly_{(\CC,\RR)}$ corresponds to the 1-category
of univariate polynomials $\Poly_\CC (1,1)$ in \cite{gambino2013}.
They write $\Poly_\CC$ for the bicategory of multivariate polynomials;
for objects $I$ and $J$,
$\Poly_\CC(I,J)$ consists of the 1-category of multivariate polynomials
indexed over $I$ and $J$.

\begin{remark}
  As this work is part of a formalisation effort
  \cite{hazratpour2024poly},
  we note a technical detail concerning the definition of morphisms
  between polynomial signatures.
  The object $f' \times_{B'} \phi_1$
  cannot be part of the data of a morphism,
  since this would get in the way of proving equality of morphisms,
  such as associativity of composition.
  One workaround is to require that a morphism $\phi : (f) \to (f')$
  provides {for any given pullback} $f \times_B \phi_1 \iso P$ a morphism
  $\phi^\#_P : P \to E$,
  such that for any two such pullbacks $P$ and $P'$,
  the triangle formed with the canonical isomorphism $i : P \to P'$ commutes.
  \[\phi^\#_P = \phi^\#_Q \circ i : P \to E\]
\end{remark}

\begin{definition}\label{def:vert-cart-ofs}
  We say that a morphism $\phi : (f : E \to B) \to (f' : E' \to B')$ in $\Poly_{(\CC,\RR)}$
  is \emph{vertical} when $\phi_1$ is an isomorphism;
  it is \emph{cartesian} when $\phi^\#$ is an isomorphism.
  We write $\vert$ and $\cart$ respectively for
  the classes of vertical and cartesian morphisms in $\Poly_{(\CC,\RR)}$.
\end{definition}
  A morphism
  $\phi : (f : E \to B) \to (f' : E' \to B')$
  can be factorised as a vertical morphism
  $(\id_B, \phi^\#) : (f : E \to B) \to (f \circ \phi^\# : Q \to B)$
  followed by a cartesian one
  $(\phi_1, \id_{Q}) : (f \circ \phi^\# : Q \to B) \to (f' : E' \to B')$,
  where $Q := f' \times_{B'} \phi_1$ is the pullback of $\phi_1$ along $f$.
  We note that $f \circ \phi^\#$ is $\RR$-exponentiable,
  since it is the pullback of $f'$.
  \[
  \begin{tikzcd}
    {Q} & E & B \\
    {E'} && {B'}
    \arrow["{\phi^\#}", from=1-1, to=1-2]
    \arrow[from=1-1, to=2-1]
    \arrow["f", from=1-2, to=1-3]
    \arrow["{\phi_1}", from=1-3, to=2-3]
    \arrow["{f'}"', from=2-1, to=2-3]
    \arrow["\lrcorner"{anchor=center, pos=0.125}, draw=none, from=1-1, to=2-3]
  \end{tikzcd}
  \quad = \quad
  \begin{tikzcd}
    {Q} & {Q} & E & B \\
    {Q} & {Q} && B \\
    {E'} &&& {B'}
    \arrow[equals, from=1-1, to=1-2]
    \arrow[equals, from=1-1, to=2-1]
    \arrow["\lrcorner"{anchor=center, pos=0.125}, draw=none, from=1-1, to=2-2]
    \arrow["{\phi^\#}", from=1-2, to=1-3]
    \arrow[equals, from=1-2, to=2-2]
    \arrow["\lrcorner"{anchor=center, pos=0.125}, draw=none, from=1-2, to=2-4]
    \arrow["f", from=1-3, to=1-4]
    \arrow[equals, from=1-4, to=2-4]
    \arrow[equals, from=2-1, to=2-2]
    \arrow[from=2-1, to=3-1]
    \arrow["\lrcorner"{anchor=center, pos=0.125}, draw=none, from=2-1, to=3-4]
    \arrow["{f \circ \phi^\#}"', from=2-2, to=2-4]
    \arrow["{\phi_1}", from=2-4, to=3-4]
    \arrow["{f'}"', from=3-1, to=3-4]
  \end{tikzcd}
  \]
\begin{proposition}
  The classes $(\vert,\cart)$
  form an orthogonal factorisation system on $\Poly_{(\CC,\RR)}$.
\end{proposition}
A detailed account of the factorisation system when $\CC = \Set$ and $\RR$ consists of all maps 
can be found in \cite[Section 5.5]{niu2025};
the construction here is not significantly different.

\section{Polynomial functors}\label{sec:polynomial-functor}

\begin{definition}\label{def:polynomial-functor-1}
  Let $f : E \to B$ be a polynomial signature in $(\CC,\RR)$.
  The polynomial functor $P_f : \RR(1) \to \RR(1)$
  associated to $f$ is given by the
  composition of functors $P_f = B_! \circ f_* \circ E^*$.
  \[\begin{tikzcd}
     & {\RR(1)} & {\RR(E)} & {\RR(B)} & {\RR(1)} & 
    \arrow["{E^*}", from=1-2, to=1-3]
    \arrow["{f_*}", from=1-3, to=1-4]
    \arrow["{B_!}", from=1-4, to=1-5]
  \end{tikzcd}\]
\end{definition}

Unlike in \cite{hua2026},
we do not require that all objects are $\RR$-objects,
and we do not require that all $\RR$-maps
are $\RR$-exponentiable.
This general setting permits us to apply this theory
to the examples considered in \cite{barton2026},
such as when $\CC = \Top$ and $\RR$ is the class of open embeddings,
or when $\CC = \Cat$ and $\RR$ is the class of Grothendieck opfibrations.

Of course, we could also define polynomial functors for
multivariate polynomial signatures.
\[\begin{tikzcd}
	I & E & B & J
	\arrow["i"', from=1-2, to=1-1]
	\arrow["f", from=1-2, to=1-3]
	\arrow["j", from=1-3, to=1-4]
\end{tikzcd}\]
\[\begin{tikzcd}
   & {\RR(J)} & {\RR(E)} & {\RR(B)} & {\RR(I)} & 
  \arrow["{j^*}", from=1-2, to=1-3]
  \arrow["{f_*}", from=1-3, to=1-4]
  \arrow["{i_!}", from=1-4, to=1-5]
\end{tikzcd}\]

\begin{remark}
  Our setting is fairly minimal for a well-behaved theory of polynomials.
  The definition of a polynomial functor already uses two of
  the three preclan axioms -- pullback stability for $E^*$ and
  closure under composition $B_!$.
  The theory of exponentiability
  (which uses the remaining preclan axiom)
  is required for $f_*$.

  Of course,
  one can also wonder if all the conditions
  for exponentiability are necessary.
  Indeed,
  we will see that to define polynomial composition,
  we would like $f$ to belong to a pullback stable class of maps $\HH$.
  We will also see that we need $f_*$ to be a partial right adjoint
  for the universal property of polynomial functors,
  and this fact should also hold for its pullbacks.
  By \Cref{lem:stable-class-of-exponentiable-maps}
  this means that $\HH \subseteq \Exp_\RR$.

  For an even more general approach,
  which may be useful for working with examples considered in \cite{anel2024}
  (in particular the example of (op)fibrations in $\Cat$ used in \Cref{sec:cat-as-types}),
  it appears that some of this theory could be recovered
  when $\RR$ is a comprehension category over $\CC$ that is not full
  (fullness is implicit in our assumptions).
  We leave this for future work.
\end{remark}

\begin{lemma}\label{lem:polynomial-functor-preserves-maps}
  Let $f : E \to B$ be a signature in $(\CC,\RR)$.
  Suppose $(\CC,\SS)$ is another preclan such that $\SS \subseteq \RR$
  and $f$ is also $\SS$-exponentiable.
  Then the polynomial functor
  \[ P_f : (\RR(1), \SS) \to (\RR(1), \SS)\]
  is a preclan morphism with respect to $(\RR(1),\SS)$.
\end{lemma}
\begin{proof}
  The polynomial functor is a composition of preclan morphisms.
  \[\begin{tikzcd}
    {(\RR(1),\SS)} & {(\RR(E),\SS)} & {(\RR(B),\SS)} & {(\RR(1),\SS)}
    \arrow["{{{E^*}}}", from=1-1, to=1-2]
    \arrow["{{{f_*}}}", from=1-2, to=1-3]
    \arrow["{{{B_!}}}", from=1-3, to=1-4]
  \end{tikzcd}\]
  That $E^* : (\RR(1),\SS) \to (\RR(E),\SS)$ is a preclan morphism is clear.
  We have already shown that $f_* : (\RR(E),\SS) \to (\RR(B),\SS)$
  is a preclan morphism in \Cref{lem:pushforward-preserves-maps}.
  The functor $B_! : \RR(B) \to \RR(1)$ clearly preserves $\SS$-maps;
  $B_!$ preserves pullbacks of $\SS$-maps since
  the inclusions $\RR(X) \subseteq \CC / X$ reflect pullbacks,
  and the map $B_! : \CC / B \to \CC / 1$ preserves pullbacks.
\end{proof}

\begin{proposition}\label{prop:polynomial-functor-preclan-morphism}
  Let $f : E \to B$ be a signature in $(\CC,\RR)$.
  Then the polynomial functor $P_f$ is a preclan morphism.
  \[ P_f : (\RR(1), \RR) \to (\RR(1), \RR)\]
  Moreover, it is a restriction of the usual polynomial functor in presheaves,
  which is also a preclan morphism.
  \[ \begin{tikzcd}
      {(\RR(1), \RR)} & {(\RR(1), \RR)} \\
      {(\wCC, \wRR)} & {(\wCC, {\wRR})}
      \arrow["{P_f}", from=1-1, to=1-2]
      \arrow["y"', hook, from=1-1, to=2-1]
      \arrow["y", hook, from=1-2, to=2-2]
      \arrow["{P_f}"', from=2-1, to=2-2]
    \end{tikzcd} \]
\end{proposition}
\begin{proof}
  This follows from \Cref{lem:polynomial-functor-preserves-maps}.
  % The polynomial functors are compositions of preclan morphisms.
  % \[\begin{tikzcd}
  %   {(\RR(1),\RR)} & {(\RR(E),\RR)} & {(\RR(B),\RR)} & {(\RR(1),\RR)} \\
  %   (\wCC,\wRR) & {(\wCC / E,\wRR)} & {(\wCC / B,\wRR)} & (\wCC,\wRR)
  %   \arrow["{{E^*}}", from=1-1, to=1-2]
  %   \arrow[hook, from=1-1, to=2-1]
  %   \arrow["{{f_*}}", from=1-2, to=1-3]
  %   \arrow[hook, from=1-2, to=2-2]
  %   \arrow["{{B_!}}", from=1-3, to=1-4]
  %   \arrow[hook, from=1-3, to=2-3]
  %   \arrow[hook, from=1-4, to=2-4]
  %   \arrow["{{E^*}}"', from=2-1, to=2-2]
  %   \arrow["{{\Pi_f}}"', from=2-2, to=2-3]
  %   \arrow["{{\Sigma_B}}"', from=2-3, to=2-4]
  % \end{tikzcd}\]
  % That the two $E^*$ are preclan morphisms is clear.
  % It follows from the equivalent conditions of \Cref{theorem:exponentiable}
  % that $f_*$ and $\Pi_f$ are preclan morphisms.
  % The map $B_!$ clearly preserves $\RR$-maps.
  % That $B_!$ preserves pullbacks of $\RR$-maps follows from the
  % observation that the pullback of an $\RR$-map
  % $f : A \to B$ along an arbitrary map $b : B' \to B$ in $\RR(X)$
  % is the same thing as the pullback of
  % $f : X_! A \to X_! B$ along $b : X_! B' \to X_! B$ in $\CC$.
  % It is straightforward to show that $\Sigma_B$ is a preclan morphism.
\end{proof}

% The previous proposition says that pullbacks of $\RR$-maps
% are preserved by polynomials.
% We now show that all connected limits are preserved when $\RR$ is a clan.

% \begin{proposition}
%   Suppose $(\CC,\RR)$ is a \emph{clan},
%   so that $\RR(1) = \CC$.
%   Let $f : E \to B$ be a signature in $(\CC,\RR)$.
%   Then $P_f : \CC \to \CC$ preserves connected limits.
% \end{proposition}
% \begin{proof}
%   By \cite[Proposition 1.16]{gambino2013},
%   we know that the polynomial functor $\widehat{P}_f : \wCC \to \wCC$
%   in presheaves associated with the image of $f$ under the Yoneda embedding
%   preserves connected limits.
%   The two functors $P_f$ and $\widehat{P}_f$
%   are related by the natural isomorphisms
%   \[
%   \begin{tikzcd}
%     \wCC & {\wCC / E} & {\wCC / B} & \wCC \\
%     \CC & {\RR(E)} & {\RR(B)} & \CC
%     \arrow["{E^*}", from=1-1, to=1-2]
%     \arrow["{f_*}", from=1-2, to=1-3]
%     \arrow["{B_!}", from=1-3, to=1-4]
%     \arrow["y", from=2-1, to=1-1]
%     \arrow["{E^*}"', from=2-1, to=2-2]
%     \arrow[from=2-2, to=1-2]
%     \arrow["{f_*}"', from=2-2, to=2-3]
%     \arrow[from=2-3, to=1-3]
%     \arrow["{B_!}"', from=2-3, to=2-4]
%     \arrow["y", from=2-4, to=1-4]
%   \end{tikzcd}
%   \]
%   The Yoneda embedding $y : \CC \to \wCC$ preserves all limits
%   and $\widehat{P}_f$ preserves connected limits.
%   Hence $y \circ P_f \iso \widehat{P}_f \circ y$ preserves connected limits.
%   Since $y : \CC \to \wCC$ is fully faithful,
%   it also reflects limits,
%   so we can conclude that $P_f$ also preserves connected limits.
%   % Note that the fully faithful functor $\RR(E) \to \wCC / E$ reflects,
%   % but does not necessarily preserve limits.
% \end{proof}

The polynomial functor associated to a signature $f : E \to B$
is characterised by its universal property:
\[ P_f X = \sum_{b : B} X^{E_b}\]

More precisely, we have the following proposition.

\begin{proposition}[Universal property of polynomial functors]\label{prop:poly-up}
  Given a signature $f : E \to B$,
  there is a bijection
  \[ \CC(\Gamma, P_f X) \iso
    \coprod_{t_1 \in \CC(\Gamma,B)}\CC(t_1 \times_B f, X)\]
  that is natural in both $\Gamma \in \CC^\op$ and $X \in \RR(1)$.
  For a given $t : \Gamma \to P_f X$,
  an element of the dependent sum can be visualised as
\[\begin{tikzcd}
	X & {t_1 \times_B f} & E \\
	& \Gamma & B
	\arrow["t_2"', from=1-2, to=1-1]
	\arrow[from=1-2, to=1-3]
	\arrow[from=1-2, to=2-2]
	\arrow["\lrcorner"{anchor=center, pos=0.125}, draw=none, from=1-2, to=2-3]
	\arrow["f", from=1-3, to=2-3]
	\arrow["t_1"', from=2-2, to=2-3]
\end{tikzcd}\]
  Of particular interest is the natural bijection applied to
  the identity $\id \in \CC(P_f X, P_f X)$,
  which results in a {universal} pair
  $\fstProj = (\id)_1 : P_f X \to B$
  and $\sndProj = (\id)_2 : \fstProj \times_B f \to X$,
  through which all others factor
\[\begin{tikzcd}[column sep = large]
	& X \\
	{t_1 \times_B f} & {\fstProj \times_B f} & E \\
	\Gamma & {P_f X} & B
	\arrow["{t_2}", from=2-1, to=1-2]
  \arrow["{t \times_B f}"', from=2-1, to=2-2]
	\arrow[from=2-1, to=3-1]
	\arrow["\lrcorner"{anchor=center, pos=0.125}, draw=none, from=2-1, to=3-2]
	\arrow["{\sndProj}"', from=2-2, to=1-2]
	\arrow[from=2-2, to=2-3]
	\arrow[from=2-2, to=3-2]
	\arrow["\lrcorner"{anchor=center, pos=0.125}, draw=none, from=2-2, to=3-3]
	\arrow["f", from=2-3, to=3-3]
	\arrow["t"', from=3-1, to=3-2]
	\arrow["{t_1}"', bend right, from=3-1, to=3-3]
	\arrow["{\fstProj}"', from=3-2, to=3-3]
\end{tikzcd}\]
\end{proposition}
\begin{proof}
  See \cite[Proposition 3.6]{hua2026}.
\end{proof}

\begin{definition}\label{def:pcomp}
  If $f : E \to B$ and $f' : E' \to B'$ are both signatures,
  then we can form the \emph{polynomial composition}
  $f \pcomp f' : \sndProj \times_{B'} f' \to P_f B'$ as follows
  \[ \begin{tikzcd}
    {E'} & {\sndProj \times_{B'} f'} & \\
    {B'} & {\fstProj \times_B f} & E \\
    & {P_fB'} & B
    \arrow["{{f'}}"', from=1-1, to=2-1]
    \arrow[from=1-2, to=1-1]
    \arrow["\lrcorner"{anchor=center, pos=0.125, rotate=-90}, draw=none, from=1-2, to=2-1]
    \arrow[from=1-2, to=2-2]
    \arrow["{f \pcomp f'}"{description, pos=0.2}, shift left=10, bend left = 50, from=1-2, to=3-2]
    \arrow["{{\sndProj}}", from=2-2, to=2-1]
    \arrow[from=2-2, to=2-3]
    \arrow[from=2-2, to=3-2]
    \arrow["\lrcorner"{anchor=center, pos=0.125}, draw=none, from=2-2, to=3-3]
    \arrow["f", from=2-3, to=3-3]
    \arrow["{{\fstProj}}"', from=3-2, to=3-3]
  \end{tikzcd}\]
  The map $f \pcomp f'$ is a polynomial signature
  since $P_f B'$ is an $\RR$-object
  and $\Exp_\RR$ is stable under pullback and closed under composition
  (\Cref{prop:exp-clan}).
\end{definition}

\begin{lemma}
  The polynomial composition ${f \pcomp f'}$
  is characterised by the following natural isomorphism.
  \[ P_{f \pcomp f'} \iso P_{f} \circ P_{f'} : \RR(1) \to \RR(1) \]
\end{lemma}
\begin{proof}
  The proof of \cite[Proposition 1.12]{gambino2013}
  works mutatis mutandis in this setting,
  where slices $\CC / X$ are replaced with $\RR(X)$.
  Indeed, we have Beck--Chevalley isomorphisms for pushforward
  (\Cref{def:beck-chevalley}),
  Beck--Chevalley isomorphisms for composition
  (\Cref{prop:beck-chevalley-left}),
  the distributivity law (\Cref{lem:distributivity}),
  and pseudofunctoriality of pullback and its adjoints. 
  % We replicate the proof of \cite[Proposition 1.12]{gambino2013}.
  % We introduce names for the maps in the above diagram
  % \[ \begin{tikzcd}
  %   {E'} & {Q'} & \\
  %   {B'} & {Q} & E \\
  %   & {P_fB'} & B
  %   \arrow["{{f'}}"', from=1-1, to=2-1]
  %   \arrow[from=1-2, to=1-1]
  %   \arrow["\lrcorner"{anchor=center, pos=0.125, rotate=-90}, draw=none, from=1-2, to=2-1]
  %   \arrow[from=1-2, to=2-2]
  %   \arrow["{p'}", from=1-2, to=2-2]
  %   \arrow["{s'}"', from=1-2, to=1-1]
  %   \arrow["{p}", from=2-2, to=3-2]
  %   \arrow["{s}", from=2-2, to=2-3]
  %   \arrow["{{\sndProj}}", from=2-2, to=2-1]
  %   \arrow["\lrcorner"{anchor=center, pos=0.125}, draw=none, from=2-2, to=3-3]
  %   \arrow["f", from=2-3, to=3-3]
  %   \arrow["{{\fstProj}}"', from=3-2, to=3-3]
  % \end{tikzcd}\]
  % Applying \Cref{lem:distributivity} to $f : E \to B$ and $E^* B'$ in $\RR(E)$,
  % we also obtain
  % \[\begin{tikzcd}
  %   {B'} \times E & {Q} & {P_f B'} \\
  %   & E & B
  %   \arrow["{E^* B'}"', from=1-1, to=2-2]
  %   \arrow["{\epsilon}"', from=1-2, to=1-1]
  %   \arrow["{p}", from=1-2, to=1-3]
  %   \arrow["{s}", from=1-2, to=2-2]
  %   \arrow["\lrcorner"{anchor=center, pos=0.125}, draw=none, from=1-2, to=2-3]
  %   \arrow["{\fstProj}", from=1-3, to=2-3]
  %   \arrow["f"', from=2-2, to=2-3]
  % \end{tikzcd}
  % \quad 
  % f_* \circ (E^* B')_! \iso \fstProj_! \circ p_* \circ \epsilon^*
  % \]
  
  % \begin{align*}
  %   & P_{f \pcomp f'} \\
  %   = \, & (P_f B')_! \circ (p \circ p')_* \circ Q'^* \\
  %   \iso \, & B_! \circ \fstProj_! \circ p_* \circ p'_* \circ s'^* \circ E'^* \\
  %   \iso \, & B_! \circ \fstProj_! \circ p_* \circ \sndProj^* \circ f'_* \circ E'^* \\
  %   \iso \, & B_! \circ \fstProj_! \circ p_* \circ \epsilon^* \circ (E^* B')^* \circ f'_* \circ E'^* \\
  %   = \, & A \\
  % \end{align*}
\end{proof}

\begin{definition}\label{def:cart-trans}
  For any cartesian morphism of signatures
  $\psi : (f : E \to B) \to (f' : E' \to B')$
  we define an associated natural transformation of polynomial functors
  $P_\psi : P_f \to P_{f'}$.

  Let us rename $\psi_1 : B \to B'$ as $\delta$.
  Since $\psi$ is cartesian, we have a pullback square
  \[ \begin{tikzcd}
    E & B \\
    {E'} & {B'}
    \arrow["f", from=1-1, to=1-2]
    \arrow["\phi"', from=1-1, to=2-1]
    \arrow["\lrcorner"{anchor=center, pos=0.125}, draw=none, from=1-1, to=2-2]
    \arrow["\delta", from=1-2, to=2-2]
    \arrow["{{f'}}"', from=2-1, to=2-2]
  \end{tikzcd}\]
  We define $P_\psi : P_f \to P_{f'}$ by whiskering the following natural transformations.
  \[\begin{tikzcd}
    {\RR(1)} & {\RR(E)} & {\RR(B)} & {\RR(1)} \\
    {\RR(1)} & {\RR(E')} & {\RR(B')} & {\RR(1)}
    \arrow["{{{E^*}}}", from=1-1, to=1-2]
    \arrow[equals, from=1-1, to=2-1]
    \arrow["\iso"{marking, allow upside down}, draw=none, from=1-1, to=2-2]
    \arrow["{{{f_*}}}", from=1-2, to=1-3]
    \arrow["\iso"{marking, allow upside down}, draw=none, from=1-2, to=2-3]
    \arrow["{{{B_!}}}", from=1-3, to=1-4]
    \arrow["\Rightarrow"{marking, allow upside down}, draw=none, from=1-3, to=2-4]
    \arrow["{{{E'^*}}}"', from=2-1, to=2-2]
    \arrow["{{{\phi^*}}}", from=2-2, to=1-2]
    \arrow["{{{f'_*}}}"', from=2-2, to=2-3]
    \arrow["{{{\delta^*}}}"', from=2-3, to=1-3]
    \arrow["{{{B'_!}}}"', from=2-3, to=2-4]
    \arrow["{{\id^*}}"', equals, from=2-4, to=1-4]
  \end{tikzcd}\]
  The first isomorphism is given by pseudofunctoriality,
  the second by Beck--Chevalley for pushforward along 
  the $\RR$-exponentiable map $f$.
  The final (non-invertible)
  natural transformation is the Beck--Chevalley
  transformation for composition applied to the commutative square
  \[ \begin{tikzcd}
    B & 1 \\
    {B'} & 1
    \arrow["{{{B}}}", from=1-1, to=1-2]
    \arrow["{{{\delta}}}", from=1-1, to=2-1]
    \arrow["{{{B'}}}"', from=2-1, to=2-2]
    \arrow[equals, from=2-2, to=1-2]
  \end{tikzcd} \]
\end{definition}

Recall that a natural transformation is cartesian
when each naturality square of the transformation is
a pullback square.
\begin{proposition}
  For any cartesian morphism of signatures
  $\psi : (f : E \to B) \to (f' : E' \to B')$
  the associated natural transformation $P_\psi : P_f \to P_{f'}$
  is cartesian.
\end{proposition}
\begin{proof}
  Unfolding the whiskering diagram defining $P_\psi$,
  we see that $P_\psi$ is a composition of three transformations,
  of which two are invertible.
  Let us call the non-invertible transformation $\alpha$.
  \[ \begin{tikzcd}
    B & 1 && {\RR(B)} & {\RR(1)} & \CC \\
    {B'} & 1 && {\RR(B')} & {\RR(1)}
    \arrow["{{{{B}}}}", from=1-1, to=1-2]
    \arrow["{{{{\delta}}}}", from=1-1, to=2-1]
    \arrow["{B_!}", from=1-4, to=1-5]
    \arrow["\alpha"{description}, Rightarrow, from=1-4, to=2-5]
    \arrow["\subseteq"{description}, draw=none, from=1-5, to=1-6]
    \arrow["{{{{B'}}}}"', from=2-1, to=2-2]
    \arrow[equals, from=2-2, to=1-2]
    \arrow["{\delta^*}", from=2-4, to=1-4]
    \arrow["{B'_!}"', from=2-4, to=2-5]
    \arrow[from=2-5, to=1-5]
  \end{tikzcd}\]
  It suffices to check that the whiskering $\alpha f_* E^*$ is cartesian,
  which is true if $\alpha$ is cartesian.
  Now, the naturality squares of the natural transformation
  $\alpha$ are always pullback squares in $\CC$,
  by the definition of Beck--Chevalley for composition.
  Since $\RR(1) \subseteq \CC$ is fully faithful,
  it reflects limits,
  hence the naturality squares are also pullback squares in $\RR(1)$.
\end{proof}

\begin{definition}\label{def:vert-trans}
  For any vertical morphism of signatures
  $\psi : (f : E \to B) \to (f' : E' \to B')$
  we define an associated natural transformation of polynomial functors
  $P_\psi : P_f \to P_{f'}$.
  Since $\psi$ is vertical,
  $\psi_1 : B \to B'$ is invertible,
  and so is its pullback.
  \[
  \begin{tikzcd}
    {f' \times_{B'} \psi_1} & E & B \\
    {E'} && {B'}
    \arrow["{{\psi^\#}}", from=1-1, to=1-2]
    \arrow["\sim"', from=1-1, to=2-1]
    \arrow["\lrcorner"{anchor=center, pos=0.0125}, draw=none, from=1-1, to=2-3]
    \arrow["f", from=1-2, to=1-3]
    \arrow["{{\psi_1}}", from=1-3, to=2-3]
    \arrow["\sim"', draw=none, from=1-3, to=2-3]
    \arrow["{{f'}}"', from=2-1, to=2-3]
  \end{tikzcd}
  \]
  Let us write simplify the diagram by writing $\delta : B' \to B$ for the inverse of $\psi_1$,
  and $\phi : E' \to E$ for the composition of $\psi^\#$ with the comparison map between $E'$ the pullback.
  \[
  \begin{tikzcd}
    {E} & {B} \\
    E' & B'
    \arrow["{f}", from=1-1, to=1-2]
    \arrow["\phi", from=2-1, to=1-1]
    \arrow["\delta"', from=2-2, to=1-2]
    \arrow["\sim"', draw=none, from=1-2, to=2-2]
    \arrow["{f'}"', from=2-1, to=2-2]
  \end{tikzcd}
  \]
  We define $P_\psi: P_f \to P_{f'}$ as the following whiskering 
  \[\begin{tikzcd}
    \RR(1) & {\RR(E)} & {\RR(B)} \\
    & {\RR(E')} & {\RR(B')} & \RR(1)
    \arrow["{E^*}", from=1-1, to=1-2]
    \arrow[""{name=0, anchor=center, inner sep=0}, "{E'^*}"', from=1-1, to=2-2]
    \arrow["{f_*}", from=1-2, to=1-3]
    \arrow["{\phi^*}", from=1-2, to=2-2]
    \arrow["\Rightarrow"{marking, allow upside down}, draw=none, from=1-3, to=2-2]
    \arrow["{\delta^*}"', from=1-3, to=2-3]
    \arrow[""{name=1, anchor=center, inner sep=0}, "{B_!}", from=1-3, to=2-4]
    \arrow["{f'_*}"', from=2-2, to=2-3]
    \arrow["{B'_!}"', from=2-3, to=2-4]
    \arrow["\iso"{marking, allow upside down}, draw=none, from=0, to=1-2]    
    \arrow["{\iso}"{marking, allow upside down}, draw=none, from=2-3, to=1]
  \end{tikzcd}\]
  where the (non-invertible) natural transformation in the middle is
  Beck--Chevalley for pushforward applied to the previous commutative square.
\end{definition}

\begin{proposition}
  For any vertical morphism of signatures
  $\psi : (f : E \to B) \to (f' : E' \to B')$,
  $P_\psi$ is vertical, meaning
  \[\begin{tikzcd}
      {P_f 1} & {P_{f'} 1} \\
      B & {B'}
      \arrow["{P_\psi 1}", from=1-1, to=1-2]
      \arrow["\sim"', draw=none, from=1-1, to=1-2]
      \arrow["\sim", from=1-1, to=2-1]
      \arrow["\fstProj"', no head, from=1-1, to=2-1]
      \arrow["\sim"', from=1-2, to=2-2]
      \arrow["\fstProj", draw=none, from=1-2, to=2-2]
      \arrow["{\psi_1}"', from=2-1, to=2-2]
      \arrow["\sim", draw=none, from=2-1, to=2-2]
    \end{tikzcd}\]
\end{proposition}

\begin{definition}
Combining the $(\vert, \cart)$ orthogonal factorisation system
with \Cref{def:cart-trans} and \Cref{def:vert-trans},
we obtain a functor from the category of polynomial signatures
to the category of endofunctors on $\RR(1)$.
\[
  P : \Poly_{(\CC,\RR)} \to [\RR(1), \RR(1)]
\]
\end{definition}

\begin{remark}
  Assuming that $\CC$ is locally Cartesian closed,
  Gambino and Kock \cite{gambino2013} use an enrichment of each slice $\CC / X$ over $\CC$
  to show that when $\CC$ is locally Cartesian closed (and taking $\RR$ to be all maps),
  the functor $P$ is faithful and the full image of $P$ consists of natural tranformations
  that are ``strong'' with respect to the given enrichment.
  When $(\CC,\RR)$ is a $\pi$-clan,
  one can produce a similar enrichment of $\RR(X)$ over $\RR(1)$,
  but when $(\CC,\RR)$ is merely a clan,
  it is unclear whether a similar argument could be made to work.
\end{remark}